\documentclass[addpoints,ngerman,12pt,answers]{exam}
\usepackage{babel}
\usepackage[top=2cm,bottom=2.5cm,left=1.5cm,right=1.5cm]{geometry}
\usepackage[utf8]{inputenc}
\usepackage[T1]{fontenc}
\usepackage{booktabs}
\usepackage{graphicx}
\usepackage{csquotes}
\usepackage{paralist}
%\usepackage{wasysym}
%\usepackage[math]{iwona}
\usepackage{setspace}
\usepackage{amsmath} 
\usepackage[math]{iwona}
\usepackage{setspace}
\usepackage{enumitem}
\usepackage{amsmath}
\usepackage{amsfonts}


\pointpoints{Point}{Points}
\bonuspointpoints{Bonuspunkt}{Bonuspunkte}
\renewcommand{\solutiontitle}{\noindent\textbf{Solution:}\enspace}

\chqword{Frage}   
\chpgword{Seite} 
\chpword{Punkte}   
\chbpword{Bonus Points} 
\chsword{Erreicht}   
\chtword{Gesamt}

\checkboxchar{\Square}
\checkedchar{\CheckedBox}

\pagestyle{headandfoot}
\runningheadrule
\firstpageheader{ECO209}{Extra Questions for Week 2}{}
\runningheader{ECO209}{Extra Questions for Week 2}{}
\firstpagefooter{ECO209}{}{\thepage\,/\,\numpages}
\runningfooter{ECO209}{}{\thepage\,/\,\numpages}
\doublespacing

\begin{document}
	\vspace*{3em}
	
	\makebox[\textwidth]{Name:\enspace\hrulefill}
	
	\vspace*{3em}
	\begin{questions}
		
		
		\question Use the table below to answer the following questions
		\begin{table}[htbp]
			\centering
			\begin{tabular}{lcccc} \hline \hline
				& $P_c$  & $Q_c$  & $P_f$  & $Q_f$ \\ \hline
				2021  & 1000  & 10    & 4     & 500 \\
				2022  & 10    & 1200  & 5     & 510 \\ \hline
			\end{tabular}%
			\label{tab:addlabel}%
		\end{table}%
		\begin{enumerate}[label=(\alph*)]
			\item Calculate the nominal GDP for the years 2021 and 2022.
			\item Calculate the real GDP for the year 2022, using the year 2021 as the base year.
			\item Calculate the real GDP for the year 2021, using the year 2022 as the base year.
			\item Calculate the real GDP growth using the prices from the year 2021.
			\item Calculate the real GDP growth using the prices from the year 2022.
			\item What accounts for the substantial disparity between the two growth rates?
		\end{enumerate}
		\clearpage
		\begin{solution}
			\begin{enumerate}[label=(\alph*)]
				\item Nominal GDP for year 2021:
				\[GDP_{21} = 1000 \times 10 + 4 \times 500 = 12000\]
				Nominal GDP for year 2022:
				\[GDP_{22} = 10 \times 1200 + 5 \times 510 = 14550\]
				\item The real GDP for year 2022 is
				\[GDP'_{22} = 1000 \times 1200 + 4 \times 510 = 1,202,040\]
				\item The real GDP for year 2021 is
				\[GDP'_{21} = 10\times 10 + 5 \times 500 = 2600\]
				\item The growth is $g_Y = 99.17 = 9917\%$
				\item The growth is $g_Y = 4.596 = 459.6\%$
				\item 
				The reason Laspeyres (part d) and Paasche (part e) Price Indexes produce differing results is rooted in their use of the prices. In both calculations, prices are fixed, at different points in time –- one fixed at initial prices (2021), while the other is fixed at end-of-period prices (2022). Consequently, the Laspeyres Index typically underestimate the actual impact of price shifts, while the Paasche Index tends to overstate this impact.
				
				In our specific scenario, consider the case of good C, where prices increased from 1000 to 10, while simultaneously witnessing a surge in quantity from 10 to 1200, indicating increased demand due to the reduced price (law of demand). When employing the Laspeyres Index, the price remains  at 1000, guided by the assumption that consumers exhibit no changes in their behavior in response to shifting prices. 
				
				Alternatively, you can think of the price as the weight for the quantity of each good. With the Laspeyres or Paasche Index way of calculating the real GDP, the weight are indeed fixed. Therefore, in the part d), the weight on the good C is remain very high, even though the market had experienced a large downward movement along the demand curve (assume only the supply curve is moving).
			\end{enumerate}
		\end{solution}		
		\clearpage
		\question Consider in an economy with two consumption goods and two investment goods, use the table and answer the question below:
		\begin{table}[htbp]
			\centering
			\begin{tabular}{rrrrrrrrr}
				& \multicolumn{1}{l}{$P_{C1}$} & \multicolumn{1}{l}{$Q_{C1}$} & \multicolumn{1}{l}{$P_{C2}$} & \multicolumn{1}{l}{$Q_{C2}$} & \multicolumn{1}{l}{$P_{I1}$} & \multicolumn{1}{l}{$Q_{I1}$} & \multicolumn{1}{l}{$P_{I2}$} & \multicolumn{1}{l}{$Q_{I2}$} \\
				\midrule
				2012  & 11    & 110   & 18    & 50    & 22    & 11    & 52    & 11 \\
				2013  & 12    & 130   & 22    & 52    & 23    & 12    & 47    & 15 \\ \hline
			\end{tabular}%
		\end{table}%
		\begin{enumerate}[label=(\alph*)]
			\item Calculate the real consumption, real investment and real GDP, using the price in 2012.
			\item Calculate the real consumption, real investment and real GDP, using the price in 2013.
			\item Calculate the Laspeyres, Paasche, and Fisher price indices for the change in real consumption, real investment, and real GDP from 2012 to 2013.
			\item Calculate Real consumption, Real investment and Real GDP in chained prices, benchmarked to 2013.
			\item Sum the real consumption and real investment and compare it to the real GDP.
		\end{enumerate}
	\end{questions}
	\clearpage
	\begin{solution}
		\begin{enumerate}[label=(\alph*)]
			\item The real consumption would be $11\times110+18\times50 = 2110$ in year 2012 and $11\times130+18\times52=2366$ in year 2013
			
			The real investment would be $22\times 11 + 52\times 11 = 814$ in year 2012 and $22\times 12 + 52\times 15 = 1044$ in year 2013
			
			The real GDP are 2924 and 3410
			\item The real consumption would be $12\times110+22\times50 = 2420$ in year 2012 and $12\times130+22\times52=2704$ in year 2013
			
			The real investment would be $23\times 11 + 47\times 11 = 770$ in year 2012 and $23\times 12 + 47\times 15 = 981$ in year 2013
			
			The real GDP are 3190 and 3685
			
			\item Laspeyres: $L_C=\frac{2366}{2110}=1.121$ $L_I=\frac{1044}{814}=1.283$ and $L_Y=\frac{3410}{2924}=1.166$
			
			Paasche: $Pi_C=\frac{2704}{2420}=1.117$ $Pi_I=\frac{981}{770}=1.274$ and $Pi_Y=\frac{3685}{3190}=1.155$
			
			Fisher: $F_F = \sqrt{1.121\times 1.117}=1.119$, $F_I = \sqrt{1.283\times 1.274}=1.278$, $F_Y=\sqrt{1.166\times 1.155}=1.161$
			
			\item The real consumption, real investment and real GDP in 2013 at 2013 prices are 2704, 981 and 3685. To calculate the chained weight value, we use the following formula
			
			\[x_{t-1}G_t = x_t \to x_{t-1} = \frac{x_t}{G_t}\]
			The $G_t$ here is the gross growth rate, which is the same as our index calculation as we did not multiply it by 100.
			
			\[C_{2012} = \frac{C_{2013}}{G_{2013}} = \frac{2704}{1.119} = 2415.7\]
			
			\[I_{2012} = \frac{I_{2013}}{G_{2013}} = \frac{981}{1.274} = 767\]
			
			\[Y_{2012} = \frac{Y_{2012}}{G_{2013}} = \frac{3685}{1.161} = 3174.9\]
			\item Summing the $C_{2012}$ and $I_{2012}$, we have $2415.7+767 = 3183.15$
			Compare this number to the real GDP in chained weight, $3183.15>3174.9$. It is because the price of consumption and investment grow at a different rate.
			\end{enumerate}
	\end{solution}		
%	\begin{center}
%		\combinedgradetable[h][questions]
%	\end{center}
	
\end{document}